\documentclass[10pt,a4paper]{amsart}
\usepackage[margin=19mm]{geometry}
\usepackage{amsmath,amssymb,mathtools}
\usepackage[scr=rsfs,cal=euler]{mathalfa}
\usepackage{xfrac}
\usepackage[unicode,hidelinks,bookmarks=false]{hyperref}
\newcommand{\ie}{i.e.\ }
\newcommand{\cf}{cf.\ }
\newcommand{\resp}{respectively\ }
\newcommand{\Subsection}{\subsection}
\providecommand{\nobreakdash}{\nobreak\hbox{-}}
\setlength{\emergencystretch}{2em}
\multlinegap0pt
\usepackage{xcolor}
\usepackage{amssymb}
\usepackage{enumerate}
\newtheorem{thm}{Theorem}
\newtheorem{prop}{Proposition}
\newcommand{\HH}{\mathbb{H}}
\newcommand{\p}{\partial}
\title{Sharp Minkowski-Type Inequality in Cartan-Hadamard 3-Spaces}
\author{Fang Hong}
\date{Source: arXiv:2603.19646v3 (CC BY 4.0)}
\begin{document}
\maketitle

\noindent\emph{Source and licence.} Sharp Minkowski-Type Inequality in Cartan-Hadamard 3-Spaces. Version arXiv:2603.19646v3 (CC BY 4.0), distributed by arXiv under the Creative Commons Attribution 4.0 International licence, http://creativecommons.org/licenses/by/4.0/, as recorded in the arXiv licence metadata for this exact version and shown on the abstract page https://arxiv.org/abs/2603.19646v3. Full licence notice: \texttt{licenses/arxiv-2603.19646-CC-BY-4.0.txt}.\par\smallskip

\noindent\textit{Editorial scope.} Counted unit: the article's prop prop:compare-inequality (source0.tex lines 599--612). The article's complete original proof of it is delivered with it (source0.tex lines 613--661, one \texttt{proof} environment of 49 lines). No mathematics is written, repaired or re-derived: the statement and the proof are the article's own lines, in the article's own notation, and no retained line is substituted or re-flowed.  Retained uncounted as the source-local context the proof needs: lines 249--256.  Nothing was omitted from any retained span, so the package carries no omission record.  No retained line contains a citation, so the wrapper carries no bibliography and no bibliography entry.  No retained line contains a figure environment, an image-inclusion command or drawing code, so no image is delivered and the package declares no asset dependency.  Editorial formatting only, in addition: an AMS \texttt{amsart} preamble that restates the article's own declarations and macros used by the retained lines, namely the environments \texttt{thm}, \texttt{prop} on separate counters, together with the standard packages those lines need.  The retained environments print in this document as Theorem 1, Proposition 1 in order of appearance, whereas the article prints the same items with the numbering of its own full document.  The complete native source of the article as distributed in the arXiv e-print for this exact version is bundled unchanged as \texttt{sources/196790/source0.tex}.
		\begin{thm}[B. Kleiner, 1992]\label{thm:isoperi1}
			Let $\Omega$ be a domain in a Cartan-Hadamard 3-space $N$ with curvature $K\leq a\leq 0$, such that its boundary $\Gamma$ is a smooth surface. Then
			\begin{equation}
				\label{eq:isoperi1}
				S(\Gamma) \geq \eta_{0,a}(V(\Omega)).
			\end{equation}
			Equality holds only if $\Omega$ is isometric to a ball in $\HH^3(a)$.
		\end{thm}

		\begin{prop}\label{prop:compare-inequality}
			Let $\Omega$ be a domain in $\HH^3$ such that its boundary $\Gamma$ is a smooth surface. Then
			\begin{align}
				\label{note2.prop}
				&\sqrt{16\pi S(\Gamma)+ 2 S(\Gamma)^2 + 2  \eta_{0,-1}(V(\Omega))^2} 
				\\
				\nonumber
				\geq
				&
				\sqrt{S(\Gamma)}\sqrt{S(\Gamma)+4\pi} + 4\pi \operatorname{arcsinh}\left( \sqrt{\frac{S(\Gamma)}{4\pi}} \right)
				+ 2 V(\Omega),
			\end{align}
			and equality holds only if $\Gamma$ is a sphere in $\HH^3$.
		\end{prop}

		\begin{proof}
			Define the two-variable functions
			\begin{align*}
				F_1(S,V)&:=\sqrt{16\pi S + 2S^2 + 2 \eta_{0,-1}(V)^2},
				\\
				F_2(S,V)&:=\sqrt{S}\sqrt{S+4\pi} + 4\pi \operatorname{arcsinh}\left( \sqrt{\frac{S}{4\pi}} \right) + 2V,
			\end{align*}
			where $S, V \geq 0$.
			We will show $F_1(S,V) \geq F_2(S,V)$ if $S \geq \eta_{0,-1}(V)$.
			
			We may compute
			\begin{align*}
				\p_{S}F_1(S,V) 
				&= \frac{16\pi + 4 S}{2\sqrt{16\pi S + 2S^2 + 2 \eta_{0,-1}(V)^2}}
				= \frac{8\pi + 2 S}{\sqrt{16\pi S + 2S^2 + 2 \eta_{0,-1}(V)^2}},
				\\
				\p_{S}F_2(S,V)
				&= \frac{2 S + 4\pi}{2\sqrt{S(S+4\pi)}}
				+
				4\pi \frac{1}{\sqrt{\frac{S}{4\pi}+1}} \frac{1}{\sqrt{4\pi}} \frac{1}{2 \sqrt{S}}
				= \frac{4\pi + S}{\sqrt{4\pi S + S^2}}.
			\end{align*}
			Therefore, for $S, V$ satisfying $S \geq \eta_{0,-1}(V)$, we have
			\begin{align}\label{Note2.tmp.compare}
				\p_{S}F_2(S,V) 
				= \frac{8\pi + 2S}{\sqrt{16\pi S + 4 S^2}}
				\leq \frac{8\pi + 2S}{\sqrt{16\pi S + 2S^2 + 2 \eta_{0,-1}(V)^2}} 
				= \p_{S}F_1(S,V),
			\end{align}
			hence
			\begin{align*}
				F_2(S,V) 
				= 
				&
				F_2(\eta_{0,-1}(V), V) + \int_{\eta_{0,-1}(V)}^{S} \p_{S}F_2(t,V) dt
				\\
				\leq
				&
				F_1(\eta_{0,-1}(V), V) + \int_{\eta_{0,-1}(V)}^{S} \p_{S}F_1(t,V) dt
				\\
				=
				&
				F_1(S,V).
			\end{align*}
			This is because $ F_2(\eta_{0,-1}(V), V) =F_1(\eta_{0,-1}(V), V)$. They both equal the total mean curvature of the sphere in $\HH^3$ with volume $V$.
			
			For a smooth closed surface \(\Gamma\) with enclosed domain $\Omega$, setting \(S=S(\Gamma)\) and \(V=V(\Omega)\), Theorem \ref{thm:isoperi1} gives $S \geq \eta_{0,-1}(V)$, and therefore the inequality follows.
			If equality holds, then the equality in \eqref{Note2.tmp.compare} holds, which implies $S = \eta_{0,-1}(V)$. By Theorem \ref{thm:isoperi1}, for a closed surface $\Gamma$ in $\HH^3$ with enclosed domain $\Omega$, if $S(\Gamma) = \eta_{0,-1}(V(\Omega))$, then $\Gamma$ is a sphere.
		\end{proof}

\end{document}
